\section{Related Work}

We group the related work into three strands: \emph{classical methods}, which
fix a transform and read the signal through it; \emph{self-supervised learning
for time series}, which learns representations without labels; and
\emph{disentangled representations for time series}, which attaches meaning to
the learned coordinates.

\subsection{Classical Methods}

Classical methods apply a \emph{fixed transform} to the raw signal to read out
its factors, and the dominant line has been the search for slowly varying
components, since the slowly varying part of a signal reflects the factors that
persist. The most direct approach splits the waveform itself: STL decomposes a
series into seasonal and trend components \cite{cleveland1990}, and EMD
extracts oscillatory modes empirically \cite{huang1998}. These deliver
\emph{components of the signal}, not the factors that generated it.

Another strand optimizes the slowness of the output directly. Slow Feature
Analysis (SFA) minimizes the temporal derivative of the output under unit
variance and pairwise decorrelation constraints, and, since the input is
expanded into low-order polynomials, it also extracts nonlinear features
\cite{wiskott2002}. Strengthening decorrelation into statistical independence
makes the slow components separate from one another too, combining slowness
with independent component analysis \cite{blaschke2007}. We use these two
variants as our label-free post-hoc rotation baselines. A further strand
changes the criterion: Predictable Feature Analysis (PFA) keeps SFA's
extraction scheme but selects the most predictable features instead of the
slowest, predictability being how well the next value is approximated from the
most recent $p$ values \cite{richthofer2013}. Yet another changes where the
signal is read: the Fourier transform reaches factors riding on the
\emph{modulation} rather than the values, and in rotating-machinery diagnosis a
bearing fault's characteristic frequencies appear not in the raw spectrum but
in the \emph{spectrum of the envelope} obtained through the Hilbert transform
\cite{randall2011}.

What this family has in common is that \emph{the transform is not learned}:
which statistics and bands to keep are fixed in advance by a human, and the
projection is refitted over the whole signal whenever the domain changes.

\subsection{Self-Supervised Learning for Time Series}

Time-series representation learning trains an encoder mapping a signal segment
to a vector without labels, taking its learning signal from the temporal
structure of the data---neighboring segments should be similar, or what is
observed now should predict what comes shortly after. The family splits into
three groups by the \emph{form of the target}.

The \emph{reconstruction family} is the largest and effectively the default.
Autoencoders pass a signal through a narrow bottleneck and restore it, the
remaining error itself being useful for anomaly detection, imputation, and
forecasting pretraining. Variational autoencoders add a prior over the latent
variables, and that prior becomes a \emph{handle for inducing structure in the
latent space}: a representative variant weights the KL term to push the
coordinates toward independence \cite{higgins2017}, a lineage leading to
Section~\ref{sec:disent}. Masked pretraining, which reconstructs randomly
masked patches, has since become established \cite{nie2023}.

The \emph{contrastive family} decides which pairs count as similar and pulls
and pushes the representations accordingly. CPC predicts the future
contrastively in the latent space \cite{oord2018}, and later time-series
methods refined those pairs: contrasting hierarchically over augmented context
views \cite{yue2022}, or using local smoothness to define stationary temporal
neighborhoods \cite{tonekaboni2021}. Their weakness is that what gets discarded
is hidden in the pair-selection rule.

The \emph{latent-prediction family} takes the representation itself, not the
raw signal, as the target \cite{lecun2022, assran2023}. Its time-series
variants differ in where the target is read: from masked patches
\cite{ennadir2025}, from the next latent token \cite{chemeris2026}, or from an
attention-pooled summary of the whole future interval \cite{petersen2026}.

This work belongs to the latent-prediction family but uses \emph{both} a
regression-form ($L_1$) and a contrastive-form (InfoNCE) loss, reporting the
two stems together to show that the horizon gate does not rely on a particular
loss form. All three groups \emph{learn} the representation, but the objective
does not determine what each coordinate means.

\subsection{Disentangled Representations for Time Series}
\label{sec:disent}

Work here has been most successful at separating \emph{factors constant over
an entire sequence} from the rest, because invariance transfers directly into
an objective. FHVAE separates what stays fixed within a sequence from what
changes across segments by assigning different priors \cite{hsu2017}, and the
disentangled sequential autoencoder splits the latent into static and dynamic
parts \cite{li2018}. On identifiability, exploiting non-stationarity by
learning to distinguish time segments makes a nonlinear ICA model identifiable
\emph{up to a pointwise transformation of each factor} when combined with
linear ICA \cite{hyvarinen2016}, and an approach placing a \emph{sparse} prior
on the change between adjacent observations followed \cite{klindt2021}. That
models cannot be selected without labels \cite{locatello2019} applies equally
to our hyperparameters, and we address it in the limitations.

What this family leaves empty is \emph{factors that vary slowly but are not
constant}: the static/non-static dichotomy pushes them into the ``dynamic''
side and separates them no further. None of the three strands creates
coordinate-block structure of the kind we seek; this work fills that place with
the axis of the prediction horizon.
